\subsection{Duomenų dalijimas}
\label{subsec:duomenu_dalijimas}

Duomenys buvo padalinti taikant $80{:}20{:}20$ metodą. T.y. vyko 2 dalijimai:
\begin{enumerate}
    \item Pirmasis dalijimas įvyko iš visos sutvarkytos duomenų aibės – $80\ \%$ mokymo + validavimo aibė bei $20\ \%$ testavimo aibė.
    \item Antrasis dalijimas įvyko iš mokymo + validavimo aibės – $80\ \%$ mokymo aibė bei $20\ \%$ validavimo aibė.
\end{enumerate}

Objektų pasiskirstymą per duomenų aibes galima pamatyti \ref{tab:objektu_pasiskirstymas} lentelėje.

\begin{table}[H]
    \centering
    \caption{Objektų pasiskirstymas po padalijimo į skirtingas duomenų aibes.}
    \begin{tabular}{|l|c|c|c|c|}
        \hline
        Aibė & Klasė $0$ & Klasė $1$ & Iš viso & Dalis \\
        \hline
        Mokymo & $282$ & $154$ & $436$ & $64\ \%$ \\ \hline
        Validavimo & $76$ & $34$ & $110$ & $16\ \%$ \\ \hline
        Testavimo & $86$ & $51$ & $137$ & $20\ \%$ \\ \hline
        \textbf{Viso} & $\mathbf{444}$ & $\mathbf{239}$ & $\mathbf{683}$ & $\mathbf{100\ \%}$ \\ \hline
    \end{tabular}
    \label{tab:objektu_pasiskirstymas}
\end{table}

Kiekviena duomenų aibė turi savo paskirtį:
\begin{itemize}
    \item Mokymo aibė – naudojama neurono svoriams ir poslinkiui atnaujinti mokymo metu.
    \item Validavimo aibė – naudojama neurono kokybei stebėti kiekvienos epochos pabaigoje. Padeda nustatyti geriausius parametrus bei pastebėti ar nevyksta persimokymas (angl.~\textit{Overfitting}).
    \item Testavimo aibė – naudojama galutiniam modelio įvertinimui. Kadangi šioje aibėje, jokie objektai nesutampa su mokymo ar validavimo aibės objektais, tai galima ištestuoti, kaip modelis veiktų su naujais, nematytais duomenimis.
\end{itemize}